\begin{frame}{SAC: Objectives and Optimization --- Critic}
\small
\begin{itemize}
    \item \textbf{Critic --- soft $Q$-function:} minimise the squared soft Bellman error, exactly as in DQN/DDPG. $\phi^{-}$ is the slow-moving target copy of the critic weights (Day 2)
\end{itemize}
\[
J_Q(\phi) = \mathbb{E}_{(\mathbf{s}_t, \mathbf{a}_t) \sim \mathcal{D}}\!\left[\, \tfrac{1}{2}\Big( Q_\phi(\mathbf{s}_t, \mathbf{a}_t) - \big( r(\mathbf{s}_t, \mathbf{a}_t) + \gamma\, \mathbb{E}_{\mathbf{s}_{t+1} \sim p}\!\left[ V_{\phi^{-}}(\mathbf{s}_{t+1}) \right] \big) \Big)^{2} \,\right]
\]
\[
V_{\phi^{-}}(\mathbf{s}_t) = \mathbb{E}_{\mathbf{a}_t \sim \pi_\theta}\!\big[\, Q_{\phi^{-}}(\mathbf{s}_t, \mathbf{a}_t) \; \textcolor{red}{-\; \alpha \log \pi_\theta(\mathbf{a}_t \mid \mathbf{s}_t)} \,\big]
\]
\[
\resizebox{\linewidth}{!}{$\displaystyle
\hat{\nabla}_\phi J_Q(\phi) = \nabla_\phi Q_\phi(\mathbf{s}_t, \mathbf{a}_t)\,\Big(\, Q_\phi(\mathbf{s}_t, \mathbf{a}_t) - \big(\, r(\mathbf{s}_t, \mathbf{a}_t) + \gamma\big[\, Q_{\phi^{-}}(\mathbf{s}_{t+1}, \mathbf{a}_{t+1}) \; \textcolor{red}{-\; \alpha \log\pi_\theta(\mathbf{a}_{t+1} \mid \mathbf{s}_{t+1})} \,\big]\, \big) \,\Big)
$}
\]
\vspace{-0.3em}
\begin{itemize}
    \item[$\Rightarrow$] {\footnotesize \textbf{This is the soft max sitting in the TD target.} Day 2 bootstrapped with $\max_{a'}Q(s',a')$; here the same slot holds $\mathbb{E}_{a'\sim\pi_\theta}[Q_{\phi^{-}} \textcolor{red}{- \alpha\log\pi_\theta}]$, which is the soft max evaluated with the actor. Cover the \textcolor{red}{red} and this is the DDPG critic update from Day 5, unchanged.}
\end{itemize}
\end{frame}

\begin{frame}{SAC: Objectives and Optimization --- Actor}
\small
\begin{itemize}
    \item \textbf{Actor:} the KL from the previous slide, $\displaystyle \arg\min_{\pi' \in \Pi} D_{\mathrm{KL}}\!\left( \pi' \,\big\|\, \exp(\tfrac{1}{\alpha} Q_\phi)/Z \right)$. Multiply through by $\alpha$ and drop $Z$ (constant in $\pi'$) and it becomes something you can just call \texttt{.backward()} on:
\end{itemize}
\[
J_\pi(\theta) = \mathbb{E}_{\mathbf{s}_t \sim \mathcal{D},\, \mathbf{a}_t \sim \pi_\theta}\!\big[\, \textcolor{red}{\alpha \log\pi_\theta(\mathbf{a}_t \mid \mathbf{s}_t)} \;-\; Q_\phi(\mathbf{s}_t, \mathbf{a}_t) \,\big]
\]
\vspace{-0.2em}
\begin{itemize}
    \item \textbf{In words:} \emph{climb the critic} (the $-Q_\phi$ term) \emph{but stay spread out} (the $+\alpha\log\pi_\theta$ term, which is $-\alpha\mathcal{H}$). DDPG's actor had only the first of those
    \vspace{0.4em}
    \item \textbf{Where the pieces came from:}
    \begin{itemize}
        \item the $\arg\min$ over $\Pi$ is soft policy improvement, now done by \emph{gradient descent} instead of exactly
        \item multiplying by $\alpha$ is cosmetic --- it does not move the $\arg\min$
        \item $Z(\mathbf{s}_t)$ is the intractable normaliser, and this is precisely where it disappears: it does not depend on $\pi'$, so its gradient is zero. \textbf{We never have to compute it}
    \end{itemize}
    \vspace{0.4em}
    \item[$\Rightarrow$] One snag remains before we can actually differentiate this --- next slide
\end{itemize}
\end{frame}

\begin{frame}{SAC: What the Actor Network Actually Outputs}
\small
\begin{itemize}
    \item \textbf{The actor is one network, and it does \emph{not} output an action.} Feed it a state and it returns \textbf{two vectors}, each with one number per action dimension:
    \begin{itemize}
        \item $\mu_\theta(\mathbf{s})$ --- the \textbf{mean}: where to aim
        \item $\sigma_\theta(\mathbf{s})$ --- the \textbf{spread}: how much to vary around that aim
    \end{itemize}
    That is the entire output. For a 6-joint arm, the actor emits 6 means and 6 spreads.
    \vspace{0.4em}
    \item \textbf{An action is then built from them in three steps:}
\end{itemize}
\vspace{-0.2em}
\[
\underbrace{\epsilon \sim \mathcal{N}(0, I)}_{\text{1. draw noise}}
\qquad
\underbrace{\mathbf{u} = \mu_\theta(\mathbf{s}) + \sigma_\theta(\mathbf{s}) \odot \epsilon}_{\text{2. scale and shift it}}
\qquad
\underbrace{\mathbf{a} = \tanh(\mathbf{u})}_{\text{3. squash into range}}
\]
\vspace{-0.2em}
\begin{itemize}
    \item \textbf{Why the $\tanh$?} Step 2 can return \emph{any} real number --- $+500$ is a perfectly legal Gaussian sample. A robot joint cannot accept a torque of $500$. $\tanh$ folds the whole real line into $(-1,1)$, so every action is in the actuator's range \emph{by construction}, and unlike clipping it stays differentiable everywhere
    \vspace{0.2em}
    \item \textbf{At evaluation time} drop step 1 and act with $\mathbf{a} = \tanh(\mu_\theta(\mathbf{s}))$ --- the mean, no noise
    \vspace{0.2em}
    \item[$\Rightarrow$] $\pi_\theta(\mathbf{a}\mid\mathbf{s})$ is the \emph{distribution} these three steps induce. $\mu_\theta$ and $\sigma_\theta$ are the only things the network learns
\end{itemize}
\end{frame}

\begin{frame}{Why That Form? It Is What Makes the Actor Trainable}
\small
\begin{itemize}
    \item \textbf{The snag:} the actor objective is an average over actions \emph{sampled from the very network we are differentiating}. You cannot backpropagate through a random draw --- ``sample from $\pi_\theta$'' has no derivative with respect to $\theta$
    \vspace{0.4em}
    \item \textbf{The fix is already in the three steps.} Look at where the randomness entered: step 1, as $\epsilon \sim \mathcal{N}(0,I)$ --- a \emph{fixed} noise source that does not depend on $\theta$ at all. Steps 2 and 3 are just multiply, add and $\tanh$: an ordinary differentiable function
    \begin{itemize}
        \item so the randomness is an \textbf{input}, not an operation we have to differentiate through
        \item gradients flow from the action back into $\mu_\theta$ and $\sigma_\theta$ as usual
    \end{itemize}
    \vspace{0.4em}
    \item This is the \textbf{reparameterization trick}. We give the three steps a single name:
\end{itemize}
\vspace{-0.2em}
\[
\mathbf{a}_t \;=\; f_\theta(\epsilon_t;\, \mathbf{s}_t) \;=\; \tanh\!\big(\mu_\theta(\mathbf{s}_t) + \sigma_\theta(\mathbf{s}_t)\odot\epsilon_t\big),
\qquad \epsilon_t \sim \mathcal{N}(0,I)
\]
\vspace{-0.2em}
\begin{itemize}
    \item[] {\footnotesize $f_\theta$ is not a new network --- it is just shorthand for ``run the actor, then steps 2 and 3''. Wherever you see $f_\theta(\epsilon_t;\mathbf{s}_t)$ from here on, read it as ``the action''.}
\end{itemize}
\end{frame}

\begin{frame}{The Squash Changes the Probabilities}
\small
\begin{itemize}
    \item Every SAC equation needs one number: $\log\pi_\theta(\mathbf{a}\mid\mathbf{s})$, \emph{how likely was the action we just played}. The Gaussian gives us the likelihood of $\mathbf{u}$ --- but we did not play $\mathbf{u}$, we played $\mathbf{a}=\tanh(\mathbf{u})$. \textbf{Those are two different numbers.}
    \vspace{0.5em}
    \item \textbf{Why they differ.} A density is \emph{probability per unit width}. $\tanh$ crushes wide stretches of $\mathbf{u}$ into narrow slivers of $\mathbf{a}$: the probability inside a stretch is unchanged, but the width is not --- so the density goes up.
\end{itemize}
\vspace{0.4em}
\begin{center}
\renewcommand{\arraystretch}{1.3}
\small
\begin{tabular}{ccc}
\hline
this range of $\mathbf{u}$ & becomes this range of $\mathbf{a}$ & density multiplied by \\
\hline
$[0,\,1]$ & $[0.00,\;0.76]$   & $\times 1.3$ \\
$[1,\,2]$ & $[0.76,\;0.96]$   & $\times 4.9$ \\
$[3,\,4]$ & $[0.995,\;0.999]$ & $\times 234$ \\
\hline
\end{tabular}
\end{center}
\vspace{0.5em}
\begin{itemize}
    \item[$\Rightarrow$] Whatever chance $\mathbf{u}$ had of landing in $[3,4]$, $\mathbf{a}$ has \emph{exactly} that chance of landing in a sliver $234\times$ narrower. Same probability, $234\times$ less room. \textbf{Near the edge of the action range the true density is hundreds of times the Gaussian's.}
\end{itemize}
\end{frame}

\begin{frame}{The Correction --- and Why You Cannot Skip It}
\small
\begin{itemize}
    \item \textbf{The fix.} Divide by how much the width shrank. That factor is just the slope of $\tanh$, namely $1-\tanh^2(u)$ --- and in logs, dividing becomes subtracting:
\end{itemize}
\vspace{-0.2em}
\[
\log \pi_\theta(\mathbf{a}\mid\mathbf{s}) \;=\; \underbrace{\log \mathcal{N}(\mathbf{u};\,\mu_\theta,\sigma_\theta)}_{\text{plain Gaussian --- easy}} \;-\; \underbrace{\textstyle\sum_i \log\!\big(1 - \tanh^2(u_i)\big)}_{\text{the correction}}
\]
\vspace{-0.2em}
\begin{itemize}
    \item[] {\footnotesize \emph{Check it against the table:} at $u=3.4$ the slope is $0.0044$, so the correction is $-\log(0.0044)=5.4$ --- a factor of $e^{5.4}\approx 225$.}
    \vspace{0.6em}
    \item \textbf{Why you cannot skip it.} The correction is \textbf{not a constant}: near $0$ in the middle of the range, around $5.4$ at the edges. Leave it out and two things break at once:
    \begin{itemize}
        \vspace{0.2em}
        \item saturated actions look \emph{less} likely than they really are, so the agent collects an inflated bonus $-\alpha\log\pi_\theta$ simply for slamming the actuator to its limit
        \vspace{0.2em}
        \item the reported entropy comes out too high, so SAC v2's tuner believes its target is already met and drives $\alpha$ down --- quietly killing the exploration the whole method exists for
    \end{itemize}
    \vspace{0.4em}
    \item[$\Rightarrow$] \emph{One line of code, and a classic silent bug if you forget it.}
\end{itemize}
\end{frame}

\begin{frame}{SAC: The Reparameterized Actor Objective}
\small
\begin{itemize}
    \item Substituting $\mathbf{a}_t = f_\theta(\epsilon_t;\mathbf{s}_t)$ into the actor objective, the expectation is now over \emph{noise} --- which does not depend on $\theta$ --- so we can differentiate straight through it:
\end{itemize}
\[
J_\pi(\theta) = \mathbb{E}_{\mathbf{s}_t \sim \mathcal{D},\, \epsilon_t \sim \mathcal{N}}\!\left[ \alpha \log \pi_\theta(f_\theta(\epsilon_t; \mathbf{s}_t) \mid \mathbf{s}_t) - Q_\phi(\mathbf{s}_t, f_\theta(\epsilon_t; \mathbf{s}_t)) \right]
\]
\[
\resizebox{0.96\linewidth}{!}{$\displaystyle
\hat{\nabla}_\theta J_\pi(\theta) = \nabla_\theta\, \alpha \log\pi_\theta(\mathbf{a}_t \mid \mathbf{s}_t) + \Big(\, \nabla_{\mathbf{a}_t} \alpha \log\pi_\theta(\mathbf{a}_t \mid \mathbf{s}_t) - \nabla_{\mathbf{a}_t} Q_\phi(\mathbf{s}_t, \mathbf{a}_t) \,\Big)\, \nabla_\theta f_\theta(\epsilon_t; \mathbf{s}_t)
$}
\]
\vspace{-0.2em}
\begin{itemize}
    \item[$\Rightarrow$] The gradient flows \emph{through the critic into the actor}, exactly like DDPG's $\nabla_\theta Q_\phi(s,\mu_\theta(s))$ from Day 5 --- now for a stochastic actor instead of a deterministic one
    \item[] {\footnotesize In code this is one line: sample $\epsilon$, build the action, evaluate $\alpha\log\pi_\theta - Q_\phi$, call \texttt{.backward()}.}
\end{itemize}
\end{frame}
